\begin{center}
% \large{\textbf{LEMBAR PENGESAHAN}} \\
\vspace{4mm}
\textbf{MODE KUASINORMAL DALAM RUANG-WAKTU SCHWARZSCHILD ANTI-DE SITTER} \\
\vspace{4mm}
\textbf{TUGAS AKHIR} \\
\vspace{4mm}
Diajukan untuk memenuhi salah satu syarat \\
memperoleh gelar Sarjana pada\\
Program Studi S-1 Fisika \\
Departemen Fisika \\
Fakultas Sains dan Analitika Data \\
Institut Teknologi Sepuluh Nopember\\
\vspace{4mm}
Oleh : \textbf{ILHAM RASYID MACHFUDI} \\
\vspace{1mm}
NRP. 5001221124 \\
\vspace{4mm}
\begin{center}
    Disetujui oleh Tim Penguji Tugas Akhir:
\end{center}

\vspace{4ex}

\begingroup
% Menghilangkan padding
\setlength{\tabcolsep}{0pt}

\noindent
\begin{tabularx}{\textwidth}{X l}
    % Data Pembimbing dari Lembar Orisinalitas
    1. Dr. rer. nat. Bintoro Anang Subagyo      & Pembimbing \\
    NIP. 197907192005011015                      &  \\
    &  \\
    &.........................  \\
    &  \\
    &  \\
    % Dosen Penguji dikosongkan (hanya nama label & titik-titik)
    2. Dr. Heru Sukamto & Penguji\\
    NIP. 198411092012121001 & \\
    &  \\
    & ......................... \\
    &  \\
    &  \\
    3. Prof. Drs. Suminar Pratapa, M.Sc, Ph.D. & Penguji\\
    NIP. 196602241990021001 &  \\
    &  \\
    & .........................\\
\end{tabularx}
\endgroup
  \begin{center}
    \textbf{SURABAYA\\Juli, 2026}
\end{center}
\end{center}

% \newpage
% \thispagestyle{empty}
% \vspace*{\fill}
%     \begin{center}
%     Halaman ini sengaja dikosongkan
%     \end{center}
% \vspace*{\fill}
% \pagebreak